%% 02-related_work.tex
%% Status: WRITTEN (v1) — does not depend on results.
%% Citations cross-checked against references.bib.

\section{Related Work}
\label{sec:related-work}


Verbal robot control requires people and robots to coordinate both the meaning of an instruction and the action it produces. We draw on research in language grounding, conversational conventions, spatial reference, shared control, and communication repair to examine how these forms of coordination develop through interaction.

\subsection{Grounding Language in Action}

Research on natural language interfaces has developed several approaches to connecting instructions with robot behavior. Generalized Grounding Graphs associate the compositional structure of language with objects, places, paths, and events to infer action plans \citep{tellex2011commands}. \citet{matuszek2013parse} developed a parser that translates route instructions into a procedural language representing actions and control structures. More recently, SayCan combined language model scores with value functions for available skills, allowing the system to consider both the relevance of an action and its feasibility for a particular robot \citep{ichter2023saycan}. Across these approaches, interpretation depends on the relationship between language, representations of the environment, and the actions available to the system.

Studies of command elicitation examine how people formulate the instructions that these systems must interpret. \citet{marge2010routes} collected spoken route descriptions across schematic, virtual, and physical scenes. Participants preferred metric descriptions but also used landmarks to specify movement. Because the robots remained stationary and participants described complete routes, the study primarily concerns instruction formulation. \citet{kocher2020left} examined a more interactive setting in which 38 children aged four to nine directed a virtual robot controlled by an experimenter, as well as their caregivers. Children used precise and vague distances, corrective feedback, and strategies that involved turning until told to stop. The experimenter assigned a movement amount to an initially unspecified distance expression and scaled later uses accordingly. These observations provide a close precedent for local units and ongoing correction, although the participants and virtual setting differ from those of the present study.

The SCOUT corpus extends this work to sustained dialogue during collaborative exploration. It combines multiple Wizard-of-Oz studies in which participants gave spoken instructions to a remotely located robot, with annotations connecting utterances to intentions and dialogue structure \citep{lukin2024scout}. Such corpora make it possible to examine instructions in relation to preceding exchanges and information about the environment. Our study builds on this tradition by considering movement magnitude, action timing, spatial reference, and interpretations of failure within a common account of verbal control across two physical navigation tasks.

\subsection{Local Conventions and Spatial Reference}

Grounding language in a computational representation is one part of establishing understanding between interaction partners. \citet{clark1991grounding} describe conversational grounding as a collaborative process through which participants establish understanding sufficient for their current purposes. Acknowledgments, subsequent contributions, and other evidence allow partners to assess this understanding, while the communication medium shapes the effort involved. \citet{brennan1996pacts} show how people develop provisional conceptualizations for referring to objects with particular partners. These conceptual pacts explain why the interpretation of an expression depends on a shared interaction history as well as its descriptive content.

Research in HRI documents several forms of linguistic adaptation. In a collaborative map task, \citet{brandstetter2017entrainment} found that a robot's choice of words influenced participants' later naming, including after the interaction had ended. \citet{asano2022alignment} found greater lexical alignment when students interacted individually with a teachable robot than in the portions of triadic interactions directed toward the robot. The relationship between alignment and rapport was more complex than predicted, suggesting that vocabulary convergence should be interpreted in relation to the interaction in which it occurs. A recent preprint by \citet{zhang2026alignment} organizes conceptual alignment according to its triggers and the levels of understanding involved. Its taxonomy was developed from dialogue between people performing concept-sorting tasks, so its application to physical robot control remains an analytical extension. Together, these accounts motivate studying local movement expressions as possible conventions whose shared meaning must be established through use.

Spatial reference introduces a related challenge because partners may describe the same environment from different perspectives. \citet{li2016spatial} found that written instructions for tabletop object selection frequently relied on a perspective that was not explicitly stated. Perspective and the complexity of spatial relations were associated with difficulties in interpreting those instructions. Although these studies examined written descriptions and subsequent object selection, they demonstrate the importance of establishing which frame a spatial expression invokes. Physical motion adds an opportunity to test that interpretation through its consequences. We examine how users respond to uncertainty about reference frames through clarification and exploratory movement, alongside their choices about the extent and timing of an action.

\subsection{Online Correction and Control Authority}

The meaning of a correction depends partly on what the robot is doing when it is issued. \citet{raux2010corrections} distinguished corrections of omitted actions, incorrect actions, and action degree, examining their linguistic, prosodic, and temporal characteristics in dialogue between people interacting with a simulated robot clerk. Their analysis connects corrective speech to the progress of an action. Dialogue can also determine who takes responsibility for resolving a problem. In collaborative teleoperation, \citet{fong2003questions} developed a system in which robots requested human assistance with perception and cognition while treating human attention as a limited resource. Although this system used a PDA interface, it established dialogue as part of the control loop. \citet{nikolaidis2017mutual} subsequently modeled human adaptability in shared autonomy, enabling a robot to adjust its behavior according to a user's willingness to change strategies.

Recent systems integrate language into ongoing control in different ways. LILAC uses language corrections during manipulation to update a learned control space that the user operates through a joystick \citep{cui2023lilac}. Language therefore changes how manual input guides the robot. BRIDGE allows users to modify planned position, velocity, and force through speech during physical assistance, while the robot provides assurances or asks for clarification \citep{wang2026bridge}. A study with 18 older adults demonstrated these adjustments across assistive tasks. The implementation pauses motion after a completed utterance while computing a modification, making the timing of an intervention an explicit feature of the interaction. \citet{cao2025interruption} addressed another aspect of timing by implementing conversational interruption handling that responds to the interrupter's intent. Their evaluation concerned discussion tasks and therefore provides evidence about managing conversational turns, rather than physical stopping performance.

This literature establishes correction and the allocation of control as central concerns in HRI. We examine how users express control arrangements in speech, including bounded movements, motion that continues until the user intervenes, repeated authorization to proceed, and delegated sequences of actions. These arrangements specify responsibility for continuation and termination alongside the intended movement. Relating them to magnitude and spatial reference helps identify what an interface must make explicit for verbal control to remain understandable and manageable. Comparing the efficiency and safety of these arrangements would require further evaluation.

\subsection{Repair, Transparency, and Calibration}

Repair requires an account of what went wrong. \citet{marge2015recovery} distinguish ambiguity about a referent from instructions that cannot be executed in the current environment. Their crowdsourced study found that recovery responses used relevant visual information to resolve ambiguity and proposed alternative plans when a request was infeasible. \citet{chai2014ground} examined how robot feedback supports common ground in situated dialogue. They found that minimal acceptance feedback could leave people believing that understanding had been established when mutual understanding was incomplete. These studies show why a successful acknowledgment cannot be assumed to identify either the robot's interpretation or the cause of a subsequent mismatch.

Providing access to intermediate processing can help users assess that interpretation. \citet{perlmutter2016transparency} exposed perception and inference outputs in a system for understanding situated language. In a study with 20 participants, visualizations on a screen improved selected measures of command communication, but the benefits did not reliably persist after the displays were removed. The findings distinguish support available during an interaction from lasting changes in a user's mental model. They also suggest that the value of transparency lies partly in helping people decide how to formulate or revise an instruction.

Recent work with language models offers further mechanisms for clarification and personalization. KnowNo uses calibrated prediction sets to determine when a planner should seek human help with an ambiguous choice \citep{ren2023knowno}. TidyBot generalizes object-placement preferences from user examples to support personalized tidying \citep{wu2023tidybot}. These capabilities provide precedents for adapting robot behavior to information supplied by users. Applying them to verbal motion control would additionally require mechanisms for establishing and revising movement units, spatial frames, and stopping arrangements in relation to the robot's observed behavior.

Our study connects these concerns by examining how users specify and regulate motion while reasoning about mismatches and possible repairs. The taxonomy relates choices about magnitude, timing, and spatial reference to users' interpretations of robot behavior. It motivates feedback that distinguishes what the system received, how it interpreted the instruction, and what it executed. We use the idea of a local control convention to interpret users' efforts to establish workable mappings within an interaction. Whether such conventions become stable, or whether particular calibration mechanisms improve control, remains a question for subsequent evaluation.
